\documentclass[11pt]{article}
\usepackage[margin=1in]{geometry}
\usepackage{amsmath,amssymb}
\usepackage[T1]{fontenc}
\usepackage{lmodern}
\usepackage{microtype}
\usepackage{xurl}
\usepackage{booktabs}
\usepackage{hyperref}
\setlength{\emergencystretch}{3em}

\title{Fertility Extension of the NIMBY Model: Key Equations and Equilibrium}
\author{Project 03: Fertility and Housing Supply}
\date{February 19, 2026}

\begin{document}
\maketitle

\section*{Purpose}
This note writes a minimal fertility extension of the existing NIMBY model, using only key equations. Equations inherited from the NIMBY paper are kept compact, and each new equation is flagged as a difference.

\section{Baseline Block Carried from NIMBY (unchanged)}

\textbf{(B1) Preferences over consumption and housing services:}
\begin{equation}
U = \sum_{j=0}^{J}\beta^{j-1}u(c_j,h_j) + \beta^J v(w_{J+1}),
\qquad
u(c,h)=\frac{\left(c^{1-\zeta}h^{\zeta}\right)^{1-\gamma}}{1-\gamma}.
\end{equation}

\textbf{(B2) Housing politics / voting condition (median-voter implementation):}
\begin{equation}
\frac{\partial V_j(a_t,h_t^{own},z_t)}{\partial p_t} + \sigma \ge 0,
\end{equation}
with equilibrium house price pinned down by the median voter's indifference condition.

\textbf{(B3) Rental-pricing / housing market block and household budget constraints:}
unchanged from the NIMBY baseline implementation.

\section{Fertility Extension: Core New Equations}

\subsection*{Difference 1: Household utility with children at home}
Use children-at-home, not total births:
\begin{equation}
\label{eq:flow_new}
 u_t
=
\frac{\left(c_t^{1-\zeta}h_{t,\text{eff}}^{\zeta}\right)^{1-\gamma}}{1-\gamma}
+
\chi\frac{n_{t,\text{home}}^{1-\eta}}{1-\eta},
\qquad
h_{t,\text{eff}}=\frac{h_t}{(1+\lambda n_{t,\text{home}})^{\psi}}.
\end{equation}
This is Candidate 2 with crowding and direct child utility.

\subsection*{Difference 2: Fertility choice in household problem}
At fertile ages, choose births $b_t \in \{0,1,2,\dots,\bar b\}$ (or continuous if needed):
\begin{equation}
V_t(s_t)=\max_{c_t,a_{t+1},h_{t+1},b_t}\left\{u_t + \beta E_t[V_{t+1}(s_{t+1})]\right\}
\end{equation}
subject to baseline budget/borrowing constraints plus child cost term
\begin{equation}
\label{eq:child_cost}
 c_t + q_t^h h_{t+1} + a_{t+1} + \kappa(b_t,n_{t,\text{home}})
\le y_t + (1+r_t)a_t + \text{housing cash-flow terms}.
\end{equation}

\subsection*{Difference 2a: Partner formation timing (baseline restriction)}
Partner formation is exogenous with age-specific hazard \(m_j\). The period sequence is:
\[
\text{partner state realized} \rightarrow \text{fertility choice} \rightarrow \text{consumption/housing/saving choices}.
\]
To keep identification clean in the baseline, we do \emph{not} model endogenous matching or rich dual-earner composition effects. This prevents two channels (partner-income effects and housing-crowding effects) from being mixed in the first implementation.

\subsection*{Difference 3: Children leave home dynamics}
Deterministic benchmark with leaving age $A_{leave}\in\{18,19\}$:
\begin{align}
N_{0,t+1} &= b_t, \\
N_{a+1,t+1} &= N_{a,t}, \quad a=0,\dots,A_{max}-1, \\
n_{t,\text{home}} &= \sum_{a=0}^{A_{leave}-1}N_{a,t}.
\end{align}
Optional robustness: stochastic departure probabilities (boomerang extension).

\subsection*{Difference 4: Endogenous generation size (demographic law of motion)}
Let $M_{j,t}$ be mass of households age $j$ at date $t$, and $\ell_j$ survival to next age:
\begin{align}
M_{j+1,t+1} &= \ell_j M_{j,t}, \quad j=1,\dots,J-1, \\
M_{1,t+1} &= \sum_{j\in\mathcal F} f_j\,M_{j,t},
\end{align}
where $f_j$ is average births per household of age $j$ implied by policy functions.

Aggregate fertility (model counterpart of TFR-like moment):
\begin{equation}
\bar n_t = \frac{\sum_{j\in\mathcal F} f_j M_{j,t}}{\sum_{j\in\mathcal F} M_{j,t}}.
\end{equation}

\section{Political-Economy Equilibrium with Endogenous Fertility}

Relative to baseline NIMBY, fertility changes equilibrium through two channels:
\begin{itemize}
\item \textbf{Current demand channel:} higher $n_{t,\text{home}}$ raises crowding costs and shifts tenure/housing demand.
\item \textbf{Future voting-mass channel:} $\bar n_t$ changes future cohort masses $M_{j,t+k}$, altering the future median voter and thus future endogenous supply/price outcomes.
\end{itemize}

Hence, the transition equilibrium is a fixed point over:
\begin{equation}
\{p_t,r_t,\text{policy}_t,M_t\}_{t\ge0}
\end{equation}
such that (i) household optimality holds, (ii) market clearing holds, (iii) voting equilibrium holds each $t$, and (iv) demographics evolve via endogenous fertility decisions.

\section{Convergence and Stationarity: What Changes}

\subsection*{1) Why convergence becomes harder}
In baseline NIMBY, demographics are exogenous, so the equilibrium loop is mainly over prices/policies given a fixed population structure. In this extension, demographics are endogenous because fertility choices today change the future age distribution.

This does \textbf{not} require strategic voting coalitions or rational expectations. Even with the same reduced-form expectations as baseline, convergence can still be slower because the model now solves a joint system:
\begin{itemize}
\item prices and voting outcomes given current demographics,
\item demographics next period given current fertility decisions,
\item then back to prices with the updated demographic composition.
\end{itemize}

So the numerical difficulty comes from an additional endogenous state transition (population composition), not from strategic behavior.

Key feedback loop:
\begin{align}
\text{housing price path}
&\rightarrow \text{fertility decisions}
\rightarrow \text{future age distribution} \nonumber\\
&\rightarrow \text{future voting coalitions}
\rightarrow \text{housing price path}.
\end{align}

\subsection*{2) Practical equilibrium concept}
Use a \textbf{transition equilibrium} from an initial demographic state, with demographic dynamics solved jointly with prices/voting. For long-run analysis:
\begin{itemize}
\item \textbf{Stationary normalized equilibrium:} normalize by total adult population if level population drifts.
\item \textbf{Balanced-growth demographic equilibrium:} if population growth converges to constant $g_M$, solve in per-capita terms.
\end{itemize}

\subsection*{3) Sufficient discipline for numerical convergence}
For tractable and stable computation, impose/target:
\begin{itemize}
\item bounded fertility choices $b_t\in[0,\bar b]$,
\item concavity in fertility utility ($\eta>0$) and convex child costs $\kappa(\cdot)$,
\item demographic update map with locally stable fixed point in normalized shares,
\item damping in outer fixed-point iterations over demographic paths.
\end{itemize}

\subsection*{4) Minimal algorithm (high-level)}
\begin{enumerate}
\item Guess a path for demographic shares and prices.
\item Solve household problem with fertility choice and children-at-home states.
\item Aggregate births to update $M_{t+1}$.
\item Recompute voting equilibrium and market clearing prices.
\item Iterate until both prices and demographic paths converge.
\end{enumerate}

\section{What Else Is Needed Before Coding}

\begin{enumerate}
\item \textbf{Definition of fertility (\(n\)):} keep this simple. Let \(n_t\) be births per fertile household in period \(t\), and target one empirical fertility moment first (a period TFR-style moment).
\item \textbf{Partner-formation block:} keep exogenous hazard \(m_j\) by age and realize partner status before fertility choice. Treat endogenous matching and rich two-earner channels as a later extension.
\item \textbf{Child-cost structure:} choose reduced-form cost $\kappa$ (time + goods) consistent with calibration data availability.
\item \textbf{Normalization for convergence:} check convergence in per-adult (or per-household) terms and age shares, not only in levels, so population growth does not look like non-convergence.
\end{enumerate}

\section{Version 0 Computational Algorithm (MATLAB-aligned)}
This implementation is designed to layer onto the current project-02 solver structure in
\texttt{Code/steadystate/}, especially \texttt{SolveSS.m}, \texttt{SolveSS\_iter.m},
\texttt{ClearMarkets.m}, and \texttt{Code\_transition\_matrix.m}.

\subsection*{State extension}
Extend the current household state to include children-at-home:
\[
s_t=(a_t,h_t,z_t,j_t,n_{t,\text{home}}).
\]

\subsection*{Outer fixed point objects}
At outer iteration \(q\), track:
\[
\mathcal X^{(q)}=\{p_t^{(q)},r_t^{(q)},M_t^{(q)}\}_{t=0}^{T_{sim}}.
\]

\subsection*{Algorithm steps}
\begin{enumerate}
\item \textbf{Household solve (modified \texttt{SolveSS\_iter}-style):}
Given \(\mathcal X^{(q)}\), solve the dynamic program with fertility choice and store policy functions
\(\Gamma(\mathcal X^{(q)})=\{c_t^*,a_{t+1}^*,h_{t+1}^*,b_t^*\}\), plus the voting derivative object \(dV/dp\).

\item \textbf{Forward density + births (new block):}
Propagate the cross-sectional distribution with fertility transitions and compute aggregate births
\[
B_t^{(q)}=\int b_t^*(s)\,d\mu_t^{(q)}(s).
\]

\item \textbf{Demographic update (new block):}
Update cohort masses:
\[
M_{1,t+1}^{(q+1)}=B_t^{(q)},\qquad
M_{j+1,t+1}^{(q+1)}=\ell_j M_{j,t}^{(q+1)}.
\]
Use normalized cohort shares in market-clearing and voting updates for numerical stability.

\item \textbf{Voting + prices (current \texttt{ClearMarkets} logic with new weights):}
Recompute political equilibrium and market-clearing prices using updated demographic weights.

\item \textbf{Damped update:}
\[
\mathcal X^{(q+1)}=(1-\xi)\mathcal X^{(q)}+\xi\,\widehat{\mathcal X}^{(q+1)},\qquad 0<\xi\le1.
\]

\item \textbf{Convergence criteria:}
Stop when all are below tolerance:
\[
\max_t|p_t^{(q+1)}-p_t^{(q)}|,\quad
\max_t\|M_t^{(q+1)}-M_t^{(q)}\|,\quad
\max_t|B_t^{(q+1)}-B_t^{(q)}|.
\]
\end{enumerate}

\subsection*{Suggested file plan}
\begin{enumerate}
\item Add \texttt{SolveSS\_iter\_fertility.m} (household DP + fertility control).
\item Add \texttt{UpdateDemographics\_fertility.m} (birth aggregation + cohort mass transition).
\item Add \texttt{ClearMarkets\_fertility.m} (outer fixed point over prices and demographics).
\item Keep a coarse-grid test mode first (small state space, short horizon) before full calibration runs.
\end{enumerate}

\section*{Numerical caution}
With endogenous fertility, the young-renter share can move voting outcomes strongly, which can induce oscillatory transitions. In practice, damping and population-share normalization are central for stable convergence.

\section*{Bottom Line}
The fertility extension is feasible with limited disruption to the NIMBY core if we: (i) keep the housing-voting block unchanged, (ii) add children-at-home utility/cost terms, (iii) explicitly evolve cohort masses through endogenous births, and (iv) solve equilibrium as a joint price-demography fixed point.

\end{document}
